\section*{अध्याय \ref{ch:g4:breathing} --- श्वासोच्छ्वास}

\begin{solution}{exo:g4:breathing:1}
साँस लेना: छाती फैलती है और हवा भीतर बहती है। साँस छोड़ना: छाती ढीली
पड़ती है और हवा बाहर बहती है।
\end{solution}

\begin{solution}{exo:g4:breathing:2}
नाक (या मुँह), \omterm{def:g4:breathing:organs}{श्वासनली}, दोनों शाखाएँ, \omterm{def:g4:breathing:organs}{फेफड़े}, और आख़िर में \omterm{prop:g4:journey-of-food:absorb}{रक्त}
वाहिकाओं में लिपटी लाखों नन्ही हवा-कोठरियाँ।
\end{solution}

\begin{solution}{exo:g4:breathing:3}
मिनट में लगभग $15$ बार। दौड़ में यह बढ़कर $40$ या उससे ज़्यादा हो जाती
है, और गहरी भी --- मेहनत करती \omterm{def:g3:bones-and-muscles:muscle}{पेशियों} को ऑक्सीजन तेज़ चाहिए, इसलिए
\omterm{def:g4:breathing:organs}{फेफड़े} उसे तेज़ी से भरते हैं।
\end{solution}

\begin{solution}{exo:g4:breathing:4}
वह ऑक्सीजन लेता है और कार्बन डाइऑक्साइड लौटाता है। अदला-बदली फेफड़ों की
नन्ही हवा-कोठरियों में, उनकी पतली दीवारों के पार, \omterm{prop:g4:journey-of-food:absorb}{रक्त} के साथ होती
है।
\end{solution}

\begin{solution}{exo:g4:breathing:5}
छोड़ी हुई हवा ज़्यादा गरम और ज़्यादा नम होती है (धुँधलाया आईना), उसमें
ऑक्सीजन कम और कार्बन डाइऑक्साइड ज़्यादा होती है।
\end{solution}

\begin{solution}{exo:g4:breathing:6}
क्योंकि हवा और \omterm{prop:g4:journey-of-food:absorb}{रक्त} की अदला-बदली को सतह चाहिए: लाखों कोठरियों की
दीवारें मिलकर आधे टेनिस मैदान जितनी मिलन-सतह छाती में तह करके भर देती
हैं। \omterm{def:g4:journey-of-food:tract}{छोटी आँत} अपनी तहों से \omterm{def:g4:journey-of-food:digestion}{पोषक तत्व} सोखने के लिए यही करामात करती है।
\end{solution}

\begin{solution}{exo:g4:breathing:7}
नाक भीतर आती हवा को रास्ते में ही गरम, नम और साफ़ कर देती है --- ठंडी,
सूखी, धूल भरी हवा नाज़ुक फेफड़ों तक पहले से तैयार होकर पहुँचती है।
\end{solution}

\begin{solution}{exo:g4:breathing:8}
उत्तर खुला है। आम तौर पर आराम में लगभग $15$ और उछल-कूद के बाद
$25$--$40$: मेहनत करती \omterm{def:g3:bones-and-muscles:muscle}{पेशियों} ने ज़्यादा ऑक्सीजन माँगी, इसलिए साँस
तेज़ और गहरी हो गई।
\end{solution}

\begin{solution}{exo:g4:breathing:9}
कि शरीर भोजन का भंडार तो उदारता से रखता है (घंटों या दिनों भर का) पर
ऑक्सीजन का भंडार लगभग रखता ही नहीं --- मुश्किल से एक मिनट का। इसलिए
ऑक्सीजन लगातार, साँस दर साँस पहुँचानी पड़ती है, और इसीलिए साँस कभी देर
तक नहीं रुक सकती।
\end{solution}

\begin{solution}{exo:g4:breathing:10}
पच्चीस साँस लेने वाले पूरी सुबह कमरे की हवा से ऑक्सीजन लेते और कार्बन
डाइऑक्साइड लौटाते रहे हैं: उसकी हवा में एक कम और दूसरी ज़्यादा हो गई
है। नुस्ख़ा नियम 1 है: खिड़कियाँ खोलो और कमरे में हवा आने दो।
\end{solution}

\begin{solution}{exo:g4:breathing:11}
दोनों \omterm{prop:g4:journey-of-food:absorb}{रक्त} और ऑक्सीजन के मिलने की जगह हैं: विशाल, पतली, \omterm{prop:g4:journey-of-food:absorb}{रक्त} से भरी
सतहें, जहाँ ऑक्सीजन भीतर उतरती है और कार्बन डाइऑक्साइड बाहर निकलती है
--- क्लोम पानी में घुली ऑक्सीजन लेते हैं, \omterm{def:g4:breathing:organs}{फेफड़े} हवा से। हवा में
ट्राउट की पंखदार क्लोम-सतहें बैठकर आपस में चिपक जाती हैं: वह विशाल
मिलन-सतह लगभग शून्य हो जाती है, और मिलन-सतह न हो तो ऑक्सीजन भी नहीं,
चाहे आस-पास की हवा कितनी भी भरपूर हो।
\end{solution}
